% Candidate A.138; method theorem, not qualifying original-gate movement.
\begin{saresult}{A.138}{Proved under stated hypotheses}{Fixed-format homeomorphism spaces are affine, but infinitesimal PL density fails}
Let $K$ be a closed bounded definable set in an o-minimal expansion
of a real closed field $R$. Fix finitely many graph-formula formats
and retain exactly the parameters defining homeomorphisms of $K$
with any prescribed finite list of definable marked-set or boundary
conditions. Let $\mathcal H$ be the definable quotient by equality
of the whole maps. For nonempty $K$ put
\[
 d(h,g)=\max_{x\in K}\|h(x)-g(x)\|
       +\max_{x\in K}\|h^{-1}(x)-g^{-1}(x)\|.
                                                        \tag{A.138.1}
\]
Then $(\mathcal H,d)$ is definably homeomorphic to a definable
subset of Euclidean space. For every $h_0\in\mathcal H$ and
$\epsilon>0$, there are a definable $d$-open neighbourhood $U$
of $h_0$, a $\delta>0$ with $B_d(h_0,\delta)\subseteq U$, and
a definable contraction $C:U\times I\to U$ such that
\[
 C(h,0)=h,\qquad C(h,1)=h_0,\qquad C(h_0,s)=h_0,
 \qquad U\subseteq B_d(h_0,\epsilon).                 \tag{A.138.2}
\]
Evaluation and inverse evaluation throughout this contraction are
jointly definable and continuous. Thus it is a contraction through
actual homeomorphisms satisfying the same marked conditions.
The empty-set case is trivial.

In contrast, if $p:[0,1]\to R$ is continuous and affine on at most
$m$ intervals, where $m\ge1$ is an ordinary finite integer, then
\[
 \max_{x\in[0,1]}|x^2-p(x)|\ge\frac1{8m^2}.           \tag{A.138.3}
\]
Over a non-Archimedean real closed field, a positive infinitesimal
$\epsilon$ therefore admits no finite PL $\epsilon$-approximation
of the semialgebraic homeomorphism $x\mapsto x^2$.
Nevertheless this homeomorphism is semialgebraically isotopic to
the identity relative to the two endpoints.

The contraction neighbourhood in (A.138.2) depends on the fixed
formula formats. The statement supplies neither a uniform
contraction over the union of all formats nor a higher
PL-to-semialgebraic concordance comparison.
\end{saresult}
\begin{proof}
Admissibility of a coefficient tuple is definable: graph uniqueness,
bijectivity, continuity, inverse continuity and the stated marked
conditions are first-order formulas. Equality of whole graphs is a
definable equivalence relation. The definable quotient is available
in an o-minimal expansion of a field; this is also the quotient
convention used in the original introduction of
\cite[p.~1]{Walsberg2015}. A finite union of formats is encoded by
a finite extra label. Evaluation and inverse evaluation are
definable on the quotient. Definable compactness gives the maxima
in (A.138.1), which is a definable metric.

We verify the precise hypothesis of Walsberg's affineness theorem
without using ordinary real separability. Suppose a definable
$d$-discrete subset $D\subseteq\mathcal H$ were infinite. Apply
definable Hardt triviality to the family of its compact graphs,
and refine the parameter partition into cells. The needed original
input works over arbitrary real closed fields
\cite[\S1.3 and Theorem~5.22, pp.~62--63]{Coste1999}.
On a positive-dimensional cell $A\subseteq D$ the graph
trivialization has the form
\[
 \theta:A\times G\xrightarrow{\sim}\Gamma|_A,
                                                        \tag{A.138.4}
\]
where $G$ is one compact graph fibre. Composing with the projection
to $(a,x)$ gives a continuous definable bijection
$A\times G\to A\times K$. It is a homeomorphism locally in $A$:
each point of the locally closed cell has a closed bounded
neighbourhood inside the cell, and on its product with $G$ the
continuous-bijection assertion follows from definable compactness.
These neighbourhoods cover $A$. Projecting instead to $(a,y)$
gives the same conclusion for inverse evaluation. Thus evaluation
and inverse evaluation are jointly continuous on $A\times K$,
as in the graph-cell argument of A.133.

Definable uniform continuity on each such compact parameter
neighbourhood times $K$ implies continuity of
$A\to(\mathcal H,d)$. Since $D$ is $d$-discrete, the inverse
image of a sufficiently small ball about any one point of $A$
is that singleton. It would be Euclidean-open in the
positive-dimensional cell, a contradiction. Hence every definable
$d$-discrete subset is finite.

Theorem~7.1 of \cite[p.~16]{Walsberg2015} now identifies
$(\mathcal H,d)$ definably with a Euclidean definable set.
Its hypotheses are stated for o-minimal expansions of arbitrary
ordered fields; no Archimedean or completeness assumption has
entered our verification. The source's discussion of uniform
metrics is in \S3.2, p.~7; its uniform affineness formulation is
Corollary~7.24, p.~26.

Definable triangulation, with the chosen point as a vertex,
provides a sufficiently small open star inside the prescribed
metric neighbourhood. Linear contraction of that star to its
vertex, transported through the definable triangulation and affine
model, gives (A.138.2). The neighbourhood is $d$-open, so it
contains some positive-radius $d$-ball. To check joint evaluation,
at a fixed $(g,y)$ use
\[
 \|h(x)-g(y)\|
 \le \max_{z\in K}\|h(z)-g(z)\|+\|g(x)-g(y)\|.
\]
The metric controls the first term and continuity of $g$ controls
the second. The same inequality applied to inverse maps proves
joint inverse continuity. Composing with $C$ proves the family
assertion, including exact identity on any parameter faces mapped
to $h_0$ when $h_0$ is the identity.

For the approximation bound, one affine interval $[a,b]$ of $p$
has length $b-a\ge1/m$. Write $c=(a+b)/2$ and
$e(x)=x^2-p(x)$. Since $p$ is affine on this interval,
\[
 e(a)+e(b)-2e(c)=\frac{(b-a)^2}{2}.                     \tag{A.138.5}
\]
The triangle inequality implies that at least one of
$|e(a)|,|e(b)|,|e(c)|$ is at least $(b-a)^2/8$, proving
(A.138.3). This uses only ordered-field arithmetic; allowing
vertices and slopes from the larger field does not change it.
A positive infinitesimal is smaller than $1/(8m^2)$ for every
ordinary finite $m$, so no finite PL map has the requested error.

Finally the explicit homotopy
\[
 H_s(x)=(1-s)x^2+s x,\qquad 0\le s\le1,
\]
has fixed endpoints, and for $x<y$ satisfies
\[
 H_s(y)-H_s(x)=(y-x)\bigl((1-s)(x+y)+s\bigr)>0.
\]
Each slice is an increasing semialgebraic homeomorphism of the
interval. The continuous definable bijection
$(x,s)\mapsto(H_s(x),s)$ on the closed bounded square has
continuous inverse by definable compactness. Thus this is a
genuine semialgebraic isotopy, not a counterexample to isotopy
or weak-homotopy comparison.

The exact midpoint and monotonicity identities, and finite sharp
piecewise-affine certificates, are checked in the attached
\texttt{check\_A138\_metric.py}. Those finite checks do not
purport to test all infinitesimals; the ordered-field proof above
supplies that assertion.
\end{proof}
